\subsection{CIFAR-10 6-animal task mismatch description}

In Table~\ref{tab:CIFAR10_Class_MisMatch_Description} we define which classes form the labeled and unlabeled set at each level of mismatch for the CIFAR-10 6 animal task. This exactly follows ~\citet{oliverRealisticEvaluationDeep2018} and creates \emph{more challenging} scenarios than other ``mismatch'' tasks on CIFAR-10 tried previously (for example, ~\citep{saitoOpenMatchOpensetConsistency2021} examine a case with all 10 classes in the unlabeled set).

\setlength{\tabcolsep}{.8cm}
\begin{table}[!h]
\centering
\begin{tabular}{l | r | r  }
& Labeled set & Unlabeled set
\\ \hline
$\zeta=0\%$ & {Bird, Cat, Deer, Dog, Frog, Horse} & {Deer, Dog, Frog, Horse}
\\
$\zeta=25\%$ & {Bird, Cat, Deer, Dog, Frog, Horse} & {\textbf{Airplane}, Dog, Frog, Horse}
\\
$\zeta=50\%$ & {Bird, Cat, Deer, Dog, Frog, Horse} & {\textbf{Airplane}, \textbf{Car}, Frog, Horse}
\\
$\zeta=75\%$ & {Bird, Cat, Deer, Dog, Frog, Horse} & {\textbf{Airplane}, \textbf{Car}, \textbf{Ship}, Horse}
\\
$\zeta=100\%$ & {Bird, Cat, Deer, Dog, Frog, Horse} & {\textbf{Airplane}, \textbf{Car}, \textbf{Ship}, \textbf{Truck}}
\\

\end{tabular}
    \caption{
    \textbf{Definition of labeled/unlabeled class mismatch scenario in CIFAR-10 6 animal task.} We bolded the non-animal classes in unlabeled set that are not in the labeled set.
    All included classes are represented with equal frequency.
    }%endcaption
    \label{tab:CIFAR10_Class_MisMatch_Description}
\end{table}


% \subsection{Illustration: OOD Unlabeled Data can be Helpful}
% \label{app:UnlabeledHelpfulMotivatingExample_MixMatch}

% We motivate the hypothesis that unlabeled data even from out-of-distribution (OOD) classes could be useful by an experiment testing the off-the-shelf performance of MixMatch~\citep{berthelotMixMatchHolisticApproach2019} on the CIFAR-10 6 animal task (with 400 images per class in training set).
% We compare MixMatch with and without perfect OOD filtering under three mismatch levels $\zeta=25\%, 50\% $ and $75\%$. 
% Results are shown in Fig.~\ref{fig:UnlabeledHelpfulMotivatingExample_MixMatch}.
% \emph{Counter-intuitively} we see that perfect OOD Filtering leads to clearly worse performance \textbf{at all mismatch levels}. This finding appears robust across 5 random train/test splits.
% This result suggests to us that unlabeled data, even with OOD classes, could be useful via MixMatch style augmentation.
% Note that our suggested Fix-A-Step procedure provides further safeguards via the gradient step modification, which are not used in Fig.~\ref{fig:UnlabeledHelpfulMotivatingExample_MixMatch}.

% \begin{figure*}[!h]
% \centering
% \begin{tabular}{c}
% \includegraphics[width=.6\textwidth]{figures/MixMatch_ToShowUnlabeledDataValuable.pdf}

% \end{tabular}
% \caption{
% \textbf{CIFAR-10 6-animal, 400 examples/class}. Results average across 5 train/test splits (shaded area shows standard deviation).
% %. Error bar showing the stanard deviation across the 5 split. 
% Both methods use the same hyper-parameters for fair comparison.
% }%endcaption	
% \label{fig:UnlabeledHelpfulMotivatingExample_MixMatch}
% \end{figure*}